% !TEX program = lualatex
\documentclass[lualatex,ja=standard]{bxjsarticle}

\usepackage{enumitem}  % リストのカスタマイズ

\title{リストのサンプル}
\author{LaTeX サンプル}
\date{\today}

\begin{document}

\maketitle

% -------------------------------------------------------------------
% itemize: 箇条書き（黒丸）
% -------------------------------------------------------------------
\section{itemize: 箇条書き}

\begin{itemize}
  \item 項目1
  \item 項目2
  \item 項目3
\end{itemize}

% -------------------------------------------------------------------
% enumerate: 番号付きリスト
% -------------------------------------------------------------------
\section{enumerate: 番号付きリスト}

\begin{enumerate}
  \item 最初の手順
  \item 次の手順
  \item 最後の手順
\end{enumerate}

% -------------------------------------------------------------------
% description: 説明リスト
% -------------------------------------------------------------------
\section{description: 説明リスト}

\begin{description}
  \item[LaTeX]  文書組版システム
  \item[BibTeX] 参考文献管理ツール
  \item[TikZ]   LaTeX 用図形描画パッケージ
\end{description}

% -------------------------------------------------------------------
% ネストされたリスト（最大4階層）
% -------------------------------------------------------------------
\section{ネストされたリスト}

\begin{itemize}
  \item レベル1
    \begin{itemize}
      \item レベル2
        \begin{itemize}
          \item レベル3
            \begin{itemize}
              \item レベル4（最深）
            \end{itemize}
        \end{itemize}
    \end{itemize}
\end{itemize}

% enumerate のネスト（(a), (i) 等の形式に変化）
\begin{enumerate}
  \item 第1項
    \begin{enumerate}
      \item 第1-1項
      \item 第1-2項
    \end{enumerate}
  \item 第2項
\end{enumerate}

% -------------------------------------------------------------------
% enumitem パッケージによるカスタマイズ
% -------------------------------------------------------------------
\section{リストのカスタマイズ (enumitem)}

% ラベルを変更
\begin{enumerate}[label=\Roman*.]
  \item ローマ数字 I
  \item ローマ数字 II
  \item ローマ数字 III
\end{enumerate}

\begin{enumerate}[label=(\alph*)]
  \item アルファベット a
  \item アルファベット b
  \item アルファベット c
\end{enumerate}

% 間隔を詰めたリスト
\begin{itemize}[noitemsep]
  \item 間隔なし項目1
  \item 間隔なし項目2
  \item 間隔なし項目3
\end{itemize}

% 開始番号を途中から指定
\begin{enumerate}[start=5]
  \item 5番目から開始
  \item 6番目
  \item 7番目
\end{enumerate}

\end{document}